\chapter{Conclusions}
\label{ch:conclusions}

% -----------------------------------------------------------------------
\section{Summary}
\label{sec:summary}

This thesis presented \textbf{Ai4k8s}, an agentic AI platform for explainable
predictive autoscaling in Kubernetes. Ai4k8s replaces the reactive threshold logic
of native HPA and VPA with a three-layer pipeline: a deterministic enforcement
layer, a local LLM reasoning layer, and a TOPSIS-based MCDA optimization layer.

The experimental evaluation on a muBench three-service microservice chain at
CrownLabs demonstrated:

\begin{itemize}
  \item Ai4k8s reduces normalized compute cost by \textbf{52\%} compared to
  native HPA at 48 concurrent connections (cost proxy 0.184 vs.\ 0.380), by
  holding replica counts conservatively when the forecast does not justify
  scale-up.

  \item VPA achieves the best raw latency (p95 = 1.09~s, 0\% SLA violations)
  by right-sizing CPU requests without adding replicas. VPA is the preferred
  choice when raw latency is the sole objective.

  \item Ai4k8s is the only system that produces a full, human-readable
  decision trace for every scaling action. MCDA agreement was full
  (gap = 0.0000) in all stable-load scenarios, confirming that the TOPSIS
  weights are well-calibrated.

  \item A 522~MB quantized language model (\textit{qwen3:0.6b}) running
  entirely on-premises achieves 1.7~s idle inference—sufficient for practical
  autoscaling—with no cloud API dependency at idle.

  \item The LLM cascade (local $\to$ Groq $\to$ regex fallback) provides
  graceful degradation: even under full CPU saturation, the pipeline completes
  within 2~s and produces a safe, deterministic recommendation.
\end{itemize}

% -----------------------------------------------------------------------
\section{Contributions}
\label{sec:contributions-ch6}

The principal contributions of this work are:

\begin{enumerate}
  \item \textbf{Three-layer autoscaling architecture.} A novel combination of
  deterministic safety rules, LLM reasoning, and MCDA optimization in a single
  production-grade pipeline. The architecture ensures that no layer alone can
  produce an unsafe or unchecked scaling decision.

  \item \textbf{On-premises LLM autoscaling.} A practical demonstration that
  sub-gigabyte quantized language models can drive Kubernetes autoscaling
  decisions at low latency on commodity hardware without GPU or cloud dependency.

  \item \textbf{Explainability by design.} The first Kubernetes autoscaler to
  log a structured, human-readable decision trace for every action, including
  LLM chain-of-thought reasoning and per-criterion MCDA scores.

  \item \textbf{Open-source implementation and evaluation harness.} All code,
  evaluation scripts, and results are publicly available at
  \url{https://ai4k8s.online} and \url{https://github.com/pedramnj/A14K8s},
  enabling reproducibility and extension.

  \item \textbf{Documented engineering gaps.} Six engineering challenges
  specific to integrating LLMs with Kubernetes autoscaling (cascade timeouts,
  parallel decision dispatch, cache invalidation, TOPSIS weight calibration)
  are documented with their root causes and solutions, providing a reference
  for future practitioners.
\end{enumerate}

% -----------------------------------------------------------------------
\section{Limitations}
\label{sec:limitations}

\begin{enumerate}
  \item \textbf{CPU starvation under load.} On a 4-vCPU VM fully saturated by
  the muBench workload, the local LLM consistently exceeds the soft timeout
  and triggers the Groq fallback. This is a hardware constraint: the VM cannot
  provide enough CPU slices for GGUF inference simultaneously with 12 active
  pods and HTTP load generation. An 8-vCPU instance is projected to resolve
  this based on the observed CPU starvation factor.

  \item \textbf{Single-node cluster.} All experiments ran on a k3s single-node
  cluster. Multi-node clusters introduce scheduling complexity, network latency,
  and node-level resource heterogeneity that may affect both the workload
  behaviour and Ai4k8s's decision quality.

  \item \textbf{Small workload scale.} The muBench chain consists of three
  services. Real production systems with hundreds of microservices and complex
  call graphs present challenges (cross-service dependency tracking, cascading
  scale events) not addressed in this work.

  \item \textbf{TOPSIS weight portability.} The default weights were calibrated
  on the muBench workload. Different workload characteristics (e.g., batch
  jobs vs.\ interactive APIs) may require different weight vectors. The system
  provides weight overrides per-deployment but does not learn weights from
  observed outcomes.

  \item \textbf{LLM hallucination risk.} Although the enforcement layer blocks
  physically impossible recommendations, the LLM may produce suboptimal
  recommendations under unusual metric combinations not seen during informal
  calibration. The MCDA layer provides a second check, but the combination
  is not formally verified.
\end{enumerate}

% -----------------------------------------------------------------------
\section{Future Work}
\label{sec:future}

\begin{enumerate}
  \item \textbf{Hardware scaling study.} Quantify the minimum vCPU count at
  which local inference remains within the soft timeout under production-level
  muBench load. Profile CPU consumption by Ollama thread count and GGUF
  quantization level to produce a hardware sizing guide.

  \item \textbf{Reinforcement learning for weight adaptation.} Replace the
  static TOPSIS weights with a reinforcement learning agent that adjusts
  weights based on observed outcomes (SLA violations, cost per cycle). This
  would allow Ai4k8s to self-tune for different workload patterns without
  manual recalibration.

  \item \textbf{Multi-cluster deployment.} Extend Ai4k8s to federated
  Kubernetes environments where autoscaling decisions must account for
  cross-cluster load balancing and inter-datacenter latency.

  \item \textbf{Fine-tuned autoscaling LLM.} Fine-tune a small language model
  (e.g., Qwen or LLaMA 1B) on a corpus of Kubernetes metrics and operator
  runbook decisions to improve recommendation accuracy on edge cases and reduce
  inference latency.

  \item \textbf{VPA and Ai4k8s hybrid mode.} Combine VPA right-sizing
  (minimizing throttling) with Ai4k8s predictive horizontal scaling (avoiding
  reactive provisioning lag) in a coordinated dual-layer controller. This
  would target both the latency advantage of VPA and the cost and explainability
  advantages of Ai4k8s simultaneously.

  \item \textbf{SLA-aware MCDA criteria.} Add an explicit SLA budget criterion
  to TOPSIS, parameterized by the deployment's contractual p95 target. This
  would allow the MCDA layer to automatically shift the cost/performance
  trade-off when an SLA breach risk is detected in the forecast.
\end{enumerate}
